\chapter{Backpropagation and Multilayer Networks}

\section*{The Question Left Open}

Does a second layer of weights add power the way more rope adds length, simply by being more of the same thing, or does it add power in some qualitatively different way? Chapter 2 ended with a proof, not a suggestion: no single-layer perceptron can represent XOR, because no straight line separates its positive examples from its negative ones. Adding a second layer of learned units between input and output was known, even at the time, to remove that ceiling in principle. A big enough network with a hidden layer can represent XOR, and a great deal more besides. Minsky and Papert said as much themselves in 1969. What they did not have, and what nobody else had in a practical, general form for another decade and a half, was an efficient way to train such a network.

The difficulty is not obvious until you try to state it precisely. Training a single-layer perceptron is straightforward because you always know exactly how wrong the output was, and you can nudge the one layer of weights directly in the direction that would have helped. Add a hidden layer, and the units in that hidden layer have no target of their own to be wrong about. Nobody labeled the correct hidden-layer activations. The only error you can measure sits at the very end of the network, after two layers of transformation, and the question this chapter answers is how to take that single number, the output error, and turn it into a specific, useful correction for every weight in both layers, including the ones with no direct connection to the output at all.

\section*{Credit Assignment}

This is usually called the credit assignment problem: when a system with many internal parts produces a wrong answer, which parts deserve to be blamed, and by how much. Backpropagation solves it with an application of the chain rule from calculus, propagating the output error backward through the network, one layer at a time, converting "the final answer was off by this much" into "this specific hidden unit contributed to that error in this specific way."

The mechanics are more approachable in code than in the notation. Below is a complete, minimal multilayer network: two inputs, one hidden layer of two units, one output, using the same kind of sigmoid activation that dominated the field for its first few decades.

\begin{Verbatim}[fontsize=\small, frame=single, xleftmargin=1em, xrightmargin=1em]
def sigmoid(x):
    return 1.0 / (1.0 + math.exp(-x))

def dsigmoid(y):
    return y * (1.0 - y)   # derivative, given the sigmoid's own output

class MLP:
    def __init__(self, n_in, n_hidden, n_out, lr=0.5):
        self.lr = lr
        self.w1 = [[random.uniform(-1, 1) for _ in range(n_in)] for _ in range(n_hidden)]
        self.b1 = [random.uniform(-1, 1) for _ in range(n_hidden)]
        self.w2 = [[random.uniform(-1, 1) for _ in range(n_hidden)] for _ in range(n_out)]
        self.b2 = [random.uniform(-1, 1) for _ in range(n_out)]

    def forward(self, x):
        h = [sigmoid(sum(w*xi for w, xi in zip(row, x)) + b)
             for row, b in zip(self.w1, self.b1)]
        o = [sigmoid(sum(w*hi for w, hi in zip(row, h)) + b)
             for row, b in zip(self.w2, self.b2)]
        return h, o

    def train_step(self, x, y):
        h, o = self.forward(x)

        o_err = [oi - yi for oi, yi in zip(o, y)]
        o_delta = [e * dsigmoid(oi) for e, oi in zip(o_err, o)]

        h_err = [sum(o_delta[k] * self.w2[k][j] for k in range(len(o_delta)))
                  for j in range(len(h))]
        h_delta = [e * dsigmoid(hi) for e, hi in zip(h_err, h)]

        for k in range(len(self.w2)):
            for j in range(len(self.w2[k])):
                self.w2[k][j] -= self.lr * o_delta[k] * h[j]
            self.b2[k] -= self.lr * o_delta[k]

        for j in range(len(self.w1)):
            for i in range(len(self.w1[j])):
                self.w1[j][i] -= self.lr * h_delta[j] * x[i]
            self.b1[j] -= self.lr * h_delta[j]
\end{Verbatim}

Two things are worth pointing at directly in this code, because they are the entire idea.

The line computing \texttt{o\_delta} is the only place the network's actual error appears. It is the discrepancy between what came out and what should have come out, scaled by how sensitive the sigmoid was at that point.

The line computing \texttt{h\_err} is backpropagation itself. Each hidden unit's share of the blame is a weighted sum of the output deltas, weighted by exactly the connection strengths, \texttt{self.w2}, that carried that hidden unit's activation forward to the output in the first place. A hidden unit that fed strongly into an output that was very wrong gets a large share of the blame. A hidden unit that barely influenced the output gets almost none. This is the chain rule, applied mechanically: the error flows backward along the same connections the activation flowed forward along, scaled at each step by how much influence that connection actually had.

Everything after those two lines is bookkeeping: use each delta to nudge the weights that produced it, in the direction that would have reduced the error.

\section*{Watching It Solve the Problem Chapter 2 Couldn't}

Train this exact network, two inputs, two hidden units, one output, on XOR.

\begin{Verbatim}[fontsize=\small, frame=single, xleftmargin=1em, xrightmargin=1em]
Training a 2-2-1 network on XOR
  epoch 1: total squared error = 1.08602
  epoch 10: total squared error = 1.08085
  epoch 100: total squared error = 1.07547
  epoch 500: total squared error = 0.78151
  epoch 1000: total squared error = 0.04895
  epoch 2000: total squared error = 0.00660
  epoch 4000: total squared error = 0.00224

Final predictions:
  input=[0, 0]  target=0  raw_output=0.0210  rounded=0
  input=[0, 1]  target=1  raw_output=0.9776  rounded=1
  input=[1, 0]  target=1  raw_output=0.9775  rounded=1
  input=[1, 1]  target=0  raw_output=0.0280  rounded=0
\end{Verbatim}

All four cases correct, with the raw outputs sitting confidently close to 0 or 1 rather than hovering near the decision boundary. This is the exact function that made Chapter 2's single-layer perceptron oscillate forever between two and four correct answers, never settling, because no line could separate its classes. A network with one hidden layer of just two units settles on a correct answer within a few thousand passes over four examples.

It's worth pausing on the shape of that error curve rather than just the ending. For the first several hundred epochs, barely anything happens: the error sits close to where it started, drifting slightly. Then, between epoch 500 and epoch 1000, it collapses by more than an order of magnitude. This is not a smooth, steady descent, and it's not an artifact of this particular toy example either; it's a first, small glimpse of a pattern that recurs at every scale this book will discuss, up to and including the largest language models in later chapters. A network spends a long stretch of training apparently making little progress while its weights are actually rearranging themselves through a difficult region, and then, once they cross into a region where a good solution is reachable by mostly local, incremental improvement, error drops quickly. What the hidden layer is doing during that quiet stretch is worth naming directly: it is finding a useful internal representation, a way of recoding the two raw inputs into two new intermediate quantities that happen to be linearly separable, even though the original inputs were not. Chapter 2 proved no line separates XOR in its original two-dimensional input space. Nothing in that proof said anything about whether a line might separate it in some other, learned space. Backpropagation's real contribution isn't just fitting weights; it's discovering a representation the original problem didn't hand you.

\section*{What This Resolves, and What It Doesn't}

The representational bottleneck from Chapter 2 is genuinely gone, in a strong sense. A network with even one hidden layer of sufficient width, given enough training and the right data, can approximate essentially any function of practical interest; this is a real, provable result, not marketing language, and it's usually called the universal approximation theorem. That is worth sitting with for a moment, because it means the ceiling Minsky and Papert identified was specific to the single-layer architecture they were analyzing, not to neural networks as a category.

It does not follow that the problem is solved outright, and the field did not treat it as solved, even after this technique became widely known in the mid-1980s through the work of David Rumelhart, Geoffrey Hinton, and Ronald Williams, who popularized and clearly demonstrated the method that had sat in Werbos's 1974 thesis largely unnoticed. Universal approximation says a suitable network exists. It says nothing about how large that network needs to be, how much data it takes to find good weights for it, or how long training takes as networks and datasets grow past the toy scale of this chapter's four-example XOR problem. Those turn out to be separate, harder questions, and this book will spend several later chapters, on convolutional networks, sequence models, and the deep networks that took another two decades to become practical, working through exactly those questions: not whether a big enough network can represent what you need, but how to actually find the weights, at a scale where "try every combination" and even "run gradient descent for a few thousand epochs on four examples" both stop being options.

That isn't where the book goes next, though, and it's worth saying why. Backpropagation's arrival didn't make every other approach to learning obsolete overnight; for the better part of two decades, a different family of methods, built on probability rather than gradient descent, delivered results neural networks of the era couldn't match, often the quiet, working alternative in the very winters Chapter 3 described. Understanding what those methods offered, and precisely where they ran out of road, turns out to matter for understanding what deep learning eventually had to displace. Chapter 5 picks that thread up first.
